% Part: sets-functions-relations
% Chapter: functions
% Section: composition

\documentclass[../../../include/open-logic-section]{subfiles}

\begin{document}

\olfileid{sfr}{fun}{cmp}
\olsection{Samenstelling: functies na elkaar gebruiken}

\begin{explain}
\oliflabeldef{sfr:fun:inv:sec}{In \olref[inv]{sec} zagen we dat de
inverse~$f^{-1}$ van !!a{bijection}~$f$ zelf ook een functie is. Nog
een manier om functies te gebruiken is ze samen te stellen: w}{W}e
kunnen uit twee functies, $f$ en~$g$, één nieuwe functie maken. We
gebruiken eerst $f$ en daarna~$g$. Dat werkt alleen wanneer alle
uitkomsten van de eerste functie geldige invoeren voor de tweede zijn.
Anders gezegd: het beeld van~$f$ moet een deelverzameling zijn van het
domein van~$g$.
\oliflabeldef{sfr:rel:ops:sec}{Functies op deze manier samenstellen is
vergelijkbaar met relaties samenstellen zoals beschreven in
\olref[rel][ops]{sec}.}{}

Een plaatje kan duidelijk maken wat er gebeurt. In
\olref{fig:composition} staan twee functies, $f \colon A \to B$ en
$g \colon B \to C$, en hun samenstelling~$(\comp{f}{g})$. De functie
$(\comp{f}{g}) \colon A \to C$ koppelt elk !!{element} van~$A$ aan
!!a{element} van~$C$. Zo bepalen we welk !!{element} van~$C$ bij
!!a{element} van $A$ hoort. Neem een invoer $x \in A$. Gebruik eerst
de functie $f$ voor~$x$; de uitkomst is $f(x) = y \in B$. Gebruik
daarna de functie $g$ voor~$y$; de uitkomst is $g(f(x)) = g(y) = z
\in C$.
\begin{figure}
  \olasset[2\olphotowidth]{assets/diagrams/composition.tikz}
  \caption{De samenstelling $g \circ f$: eerst $f$ en daarna~$g$.}
  \ollabel{fig:composition}
\end{figure}
\end{explain}

\begin{defn}[Samenstelling]
Neem de functies $f\colon A \to B$ en $g\colon B \to C$. De
\emph{samenstelling} van $f$ met~$g$ is $\comp{f}{g} \colon A \to C$.
Voor elke invoer geldt $(\comp{f}{g})(x) = g(f(x))$.
\end{defn}

\begin{ex}
Neem de functies $f(x) = x + 1$ en $g(x) = 2x$. Omdat
$(\comp{f}{g})(x) = g(f(x))$, tel je bij elke invoer~$x$ eerst één op.
Daarna verdubbel je de uitkomst. De samenstelling is dus
$(\comp{f}{g})(x) = 2(x+1)$.
\end{ex}

\begin{prob}
Bewijs dat als $f \colon A \to B$ en $g \colon B \to C$ allebei
!!{injective} zijn---verschillende invoeren geven dus verschillende
uitkomsten---ook $\comp{f}{g}\colon A \to C$ !!{injective} is.
\end{prob}

\begin{prob}
Bewijs dat als $f \colon A \to B$ en $g \colon B \to C$ allebei
!!{surjective} zijn---ze bereiken dus elk element van hun
doelverzameling---ook $\comp{f}{g}\colon A \to C$ !!{surjective} is.
\end{prob}

\begin{prob}
Neem $f \colon A \to B$ en $g \colon B \to C$. De grafiek van
$\comp{f}{g}$ moet $R_f \mid R_g$ zijn. Laat dat zien.
\end{prob}
\end{document}
